\documentclass[10.5pt]{article}
\usepackage[UTF8]{ctex}
\usepackage[a4paper,top=1.5cm,bottom=1.5cm,left=1.4cm,right=1.4cm]{geometry}
\usepackage{booktabs}
\usepackage{longtable}
\usepackage{array}
\usepackage{xcolor}
\definecolor{accent}{HTML}{1F4E79}
\newcommand{\sect}[1]{\par\vspace{9pt}{\large\bfseries\color{accent}#1}\par\vspace{3pt}}
\newcolumntype{L}[1]{>{\raggedright\arraybackslash}p{#1}}
\setlength{\parindent}{0pt}
\setlength{\parskip}{2pt}
\begin{document}
\begin{center}
{\LARGE\bfseries 工创小车 · 任务清单}\par\vspace{4pt}
{\small 重走古丝路 · 自动循迹避障小车（华南理工大学 · 工程训练分散项目）}\par\vspace{2pt}
{\small 5 人 · 13 项任务 · v2.3 · 更新：2026-09-20（第 3 周末）}\par
\end{center}
\vspace{-2pt}

\sect{0. 关键时间与红线}
\begin{tabular}{@{}L{2.6cm}L{14.4cm}@{}}
\toprule
\textbf{项} & \textbf{内容}\tabularnewline
\midrule
现在 & 9/20（周日）· 第 3 周末\tabularnewline
中秋节 & 9/25\textasciitilde{}9/27（人在校，照常推进）\tabularnewline
国庆放假 & 10/1\textasciitilde{}10/7（只排软件类任务）\tabularnewline
内部全功能冻结 & 10/25（周日）—— 4 障全过 + 精停 + 全套材料\tabularnewline
预测比赛周 & 10/26\textasciitilde{}11/1（比赛 + 演讲）\tabularnewline
\bottomrule
\end{tabular}

\par\vspace{4pt}
{\small \textbf{红线（违反要扣分或没资格）}：整车尺寸不超过 250×200×200mm（超标每轮扣 10\%）；全车只能一块电源（多电源扣 10\%）；车架必须自制，主结构自己设计 + 中心设备加工 + 全程留照片/视频（买现成一票否决）；全自动运行，比赛禁遥控；避障不得接触障碍物；与障碍物碰一下即该障不得分；终点停车：车尾投影不得超过终点线（20 分）}\par

\sect{1. 编号规则}
任务号 = 一个字母 + 两位数字；字母看类别，数字从 01 起。\par\vspace{3pt}
\begin{tabular}{@{}L{1.2cm}L{4.2cm}L{11.6cm}@{}}
\toprule
\textbf{字母} & \textbf{类别} & \textbf{任务}\tabularnewline
\midrule
M & 机械 & M01 结构设计与建模 · M02 加工与装配\tabularnewline
H & 硬件 & H01 选型与采购 · H02 电路与电气集成\tabularnewline
A & 算法 & A01 循迹控制 · A02 避障与回线\tabularnewline
P & 平台（软件基础） & P01 反馈链路与工具链 · P02 底层驱动与框架\tabularnewline
D & 文档与协作 & D01 项目仓库与协作规范 · D02 文档与展示材料\tabularnewline
F & 财务 & F01 账务与报销\tabularnewline
S & 辅助 & S01 辅助电路 · S02 辅助装配\tabularnewline
\bottomrule
\end{tabular}

\sect{2. 分工总览}
\begin{tabular}{@{}L{3.4cm}L{13.6cm}@{}}
\toprule
\textbf{成员} & \textbf{负责任务}\tabularnewline
\midrule
李沛阳（组长） & D01 仓库与协作规范 · P01 反馈链路与工具链 · P02 底层驱动与框架 · F01 账务与报销\tabularnewline
王政淋 & H01 选型与采购 · H02 电路与电气集成\tabularnewline
熊雅琦 & A01 循迹控制 · S01 辅助电路 · S02 辅助装配\tabularnewline
李昕程 & M01 结构设计与建模 · M02 加工与装配\tabularnewline
毕烨恺 & A02 避障与回线 · D02 文档与展示材料\tabularnewline
一起做 & 车架与 PCB 对接约定（王政淋 + 李昕程，熊雅琦辅助）\tabularnewline
\bottomrule
\end{tabular}

\sect{3. 任务详情（接下来要做的）}
{\small 每行读法：做 什么 → 交 给谁什么 → 需 要谁的什么 → 什么时候要。已完成的事在文末 §8。\par}
\vspace{5pt}{\bfseries 李沛阳（组长）—— P01 · P02 · F01}\par\vspace{2pt}
{\small
\begin{longtable}{@{}L{1.0cm}L{2.3cm}L{9.7cm}L{4.8cm}@{}}
\toprule
\textbf{编号} & \textbf{任务} & \textbf{做 · 交 · 需} & \textbf{截止日期}\tabularnewline
\midrule
\endhead
P01 & 反馈链路与工具链 & 做：接真实传感器，读数落盘出图；再做 PID/状态机数据与在线调参全链路；交：熊雅琦、毕烨恺调试时直接用的一整套工具；需：王政淋的面包板 demo（9/28）先接真实模块；李昕程的实车 v1（10/13）后跑全链路 & 阶段2 9/30；阶段3 10/14\tabularnewline
P02 & 底层驱动与框架 & 做：写电机 PWM、编码器测速、串口、I2C、超声波、电压检测；9/30 出初版，先在面包板 demo 上冒烟，10/14 冻结；交：熊雅琦、毕烨恺的驱动接口说明（函数怎么调、数据长什么样）；需：王政淋的面包板 demo（9/28）；王政淋校验后的《引脚分配表》（9/28 冻结） & 初版 9/30；冻结 10/14\tabularnewline
F01 & 账务与报销 & 做：收集发票与商品照片、填成本台账、赛后统一报账；交：毕烨恺的成本台账（进 PPT）；需：全组的电子发票 + 收货照片 + 一句用途说明（10/21 前收齐） & 收齐 10/21；赛后一周内报账\tabularnewline
\bottomrule
\end{longtable}

}

\vspace{5pt}{\bfseries 王政淋 —— H01 · H02}\par\vspace{2pt}
{\small
\begin{longtable}{@{}L{1.0cm}L{2.3cm}L{9.7cm}L{4.8cm}@{}}
\toprule
\textbf{编号} & \textbf{任务} & \textbf{做 · 交 · 需} & \textbf{截止日期}\tabularnewline
\midrule
\endhead
H01 & 选型与采购 & 做：电子元件到货清点、缺件补单；交：全组齐套可用的电子元件 & 清点 9/30\tabularnewline
H02 & 电路与电气集成 & 做：校验接线设计 → 落原理图 → 搭面包板 demo 交李沛阳（9/28）→ 边等冒烟反馈边布板 → 9/30 前下单（冒烟有问题则 10/3 前改板重下）→ 收货焊接 → 10/11 点亮并交接线图；交：李沛阳的面包板 demo（9/28）；李沛阳接线图与点亮结果（10/11）；需：李沛阳的《引脚分配表》v0.2；李昕程的关键尺寸与《车架与 PCB 对接约定》（9/24）；熊雅琦的元件与封装核对（9/28）；李沛阳的冒烟反馈（10/3 前） & 原理图+约定 9/24；面包板 demo 9/28；下单 9/30 前；点亮 10/11\tabularnewline
\bottomrule
\end{longtable}

}

\vspace{5pt}{\bfseries 熊雅琦 —— A01 · S01 · S02}\par\vspace{2pt}
{\small
\begin{longtable}{@{}L{1.0cm}L{2.3cm}L{9.7cm}L{4.8cm}@{}}
\toprule
\textbf{编号} & \textbf{任务} & \textbf{做 · 交 · 需} & \textbf{截止日期}\tabularnewline
\midrule
\endhead
A01 & 循迹控制 & 做：灰度处理 + PID，处理弯道与回环；本人上车调参，10/18 能跑弯、10/21 出全赛道稳定版；交：毕烨恺的全赛道稳定循迹版（他做避障回线的基座）；需：李沛阳的驱动接口初版（9/30，先离线写代码）；王政淋的面包板 demo（9/28）先验灰度读数；李昕程的实车 v1（10/13） & 9/30 起可编；10/14 首次上车；10/18 能跑弯；10/21 稳定版\tabularnewline
S01 & 辅助电路 & 做：打板前核对元件、封装与接线，避免打板出错；交：王政淋一份核对过的落图依据；需：王政淋的 \texttt{netlist.json} 与原理图草案 & 9/28 前\tabularnewline
S02 & 辅助装配 & 做：装配打下手、拍照记录、固定走线；交：李昕程的装配照片与记录；需：李昕程的装配安排与到货材料 & 10/8\textasciitilde{}10/13；定版 10/17\textasciitilde{}10/18\tabularnewline
\bottomrule
\end{longtable}

}

\vspace{5pt}{\bfseries 李昕程 —— M01 · M02}\par\vspace{2pt}
{\small
\begin{longtable}{@{}L{1.0cm}L{2.3cm}L{9.7cm}L{4.8cm}@{}}
\toprule
\textbf{编号} & \textbf{任务} & \textbf{做 · 交 · 需} & \textbf{截止日期}\tabularnewline
\midrule
\endhead
M01 & 结构设计与建模 & 做：整车模型、关键尺寸、关键孔位、加工图，导出件传 \texttt{mech/设计文件/}；与王政淋一起定《车架与 PCB 对接约定》；交：王政淋的板形依据、M02 的加工图、全组的尺寸确认（红线 $\leq$ 250×200×200mm）；需：王政淋的板子需求（对接讨论里定）和电控件尺寸 & 设计文件 9/24；模型冻结 9/28\tabularnewline
M02 & 加工与装配 & 做：加工材料与紧固件自己买（9/24 前下单）；用中心设备加工并全程拍照录像；装出实车 v1，10/18 定版车架；交：熊雅琦、毕烨恺的实车；毕烨恺的加工证据；需：自己的加工图（9/24）；王政淋的点亮结果（10/11）用于装电控 & 加工 10/8\textasciitilde{}10/11；实车 v1 10/13；定版 10/18\tabularnewline
\bottomrule
\end{longtable}

}

\vspace{5pt}{\bfseries 毕烨恺 —— A02 · D02}\par\vspace{2pt}
{\small
\begin{longtable}{@{}L{1.0cm}L{2.3cm}L{9.7cm}L{4.8cm}@{}}
\toprule
\textbf{编号} & \textbf{任务} & \textbf{做 · 交 · 需} & \textbf{截止日期}\tabularnewline
\midrule
\endhead
A02 & 避障与回线 & 做：超声波避障：单测 → 检测 demo → 绕行 + 回线 → 终点精停；本人上车调参；交：全组的成品功能：4 障全过 + 精停；需：李沛阳的遥测与驱动接口（9/30）；王政淋的面包板 demo（9/28）先验超声波读数；李昕程的实车 v1（10/13）；熊雅琦的稳定循迹版（10/21） & 单测 10/14；demo 10/18；绕行回线 10/21；达标 10/25\tabularnewline
D02 & 文档与展示材料 & 做：做 PPT、答辩问答、演示视频；收全员素材；交：全组一套完整的展示材料；需：全组素材（10/18 前）；李沛阳的成本台账（10/21）；李昕程的加工证据（加工期间） & 草稿 10/21；冻结 10/25\tabularnewline
\bottomrule
\end{longtable}

}

\vspace{6pt}{\bfseries 共同任务：车架与 PCB 对接约定（王政淋 + 李昕程，熊雅琦辅助）}\par\vspace{3pt}
两人一起商量，把“板子做多大、装在底板哪里、打几个孔、线从哪边出来、上面留多高”写清楚，填进 \texttt{mech/车架与PCB对接约定.md}。只写对方必须知道的尺寸，细节两人自己定，不用写很长。\par\vspace{2pt}
\textbf{交：}王政淋画 PCB 的尺寸依据、李昕程装车的同一份尺寸　\textbf{需：}熊雅琦帮忙核对元件位置\par\vspace{2pt}
\textbf{截止：}9/24 前定稿；之后要改，两人说一声即可（别单方面改）\par

\sect{4. 交接时间线（按节点）}
{\small 只列重要交接日；按周与节假日分断。\par}
\vspace{5pt}{\bfseries\color{accent}9/21\textasciitilde{}9/27 · 原理图与面包板}\par\vspace{2pt}
\begin{itemize}\setlength{\itemsep}{1pt}\setlength{\parskip}{0pt}
\item \textbf{9/24}：王政淋 ＋ 李昕程 定稿《车架与 PCB 对接约定》（熊雅琦辅助核对）
\item \textbf{9/24}：王政淋 落定原理图；李沛阳《引脚分配表》校验完成
\item \textbf{9/24}：李昕程 交设计文件（整车模型 ＋ 关键尺寸 ＋ 加工图）→ \texttt{mech/设计文件/}
\item \textbf{9/24 前}：李昕程 自行下单加工材料与紧固件
\item \textbf{9/25\textasciitilde{}9/27}：王政淋 按原理图搭面包板 demo（双 I2C、灰度、超声波、编码器、TB6612 逻辑、降压与 ADC 分压）
\item 面包板只验信号级；大电流与电源保护留到 PCB 验证
\end{itemize}
\vspace{5pt}{\bfseries\color{accent}9/28\textasciitilde{}9/30 · 关键下单窗口}\par\vspace{2pt}
\begin{itemize}\setlength{\itemsep}{1pt}\setlength{\parskip}{0pt}
\item \textbf{9/28}：王政淋 交面包板 demo → 李沛阳；李沛阳 开始桌面冒烟
\item \textbf{9/28}：《引脚分配表》冻结；李昕程 模型冻结；熊雅琦 完成元件与封装核对
\item \textbf{9/30}：王政淋 PCB 下单 ＋ 电子元件到货清点；李沛阳 交驱动接口初版
\end{itemize}
\vspace{5pt}{\bfseries\color{accent}10/1\textasciitilde{}10/7 · 国庆假（硬件停）}\par\vspace{2pt}
\begin{itemize}\setlength{\itemsep}{1pt}\setlength{\parskip}{0pt}
\item 李沛阳：驱动逐个冒烟、P02 初版修订
\item 熊雅琦、毕烨恺：离线写算法代码
\item \textbf{10/3 前}：冒烟发现原理图问题就改板重下；没问题按 9/30 下单推进
\end{itemize}
\vspace{5pt}{\bfseries\color{accent}10/8\textasciitilde{}10/11 · 节后硬件攻坚}\par\vspace{2pt}
\begin{itemize}\setlength{\itemsep}{1pt}\setlength{\parskip}{0pt}
\item \textbf{10/8\textasciitilde{}10/11}：李昕程 中心设备加工车架与底板（全程拍照录像）
\item \textbf{10/11}：王政淋 交 PCB 点亮 ＋ 接线图 → 李沛阳
\end{itemize}
\vspace{5pt}{\bfseries\color{accent}10/12\textasciitilde{}10/18 · 上车调试周}\par\vspace{2pt}
\begin{itemize}\setlength{\itemsep}{1pt}\setlength{\parskip}{0pt}
\item \textbf{10/13}：李昕程 交实车 v1（能通电、能推着走）→ 熊雅琦、毕烨恺
\item \textbf{10/14}：李沛阳 交驱动接口冻结版 ＋ 遥测上车；熊雅琦、毕烨恺开始上车调参
\item \textbf{10/18}：全组 交素材 → 毕烨恺；李昕程 交定版车架
\end{itemize}
\vspace{5pt}{\bfseries\color{accent}10/19\textasciitilde{}10/25 · 收口周}\par\vspace{2pt}
\begin{itemize}\setlength{\itemsep}{1pt}\setlength{\parskip}{0pt}
\item \textbf{10/21}：熊雅琦 交稳定循迹版（全赛道）→ 毕烨恺；李沛阳 交成本台账 → 毕烨恺
\item \textbf{10/25}：毕烨恺 交 4 障全过 ＋ 精停 ＋ 全套材料
\end{itemize}
\vspace{5pt}{\bfseries\color{accent}10/26\textasciitilde{}11/1 · 比赛周}\par\vspace{2pt}
\begin{itemize}\setlength{\itemsep}{1pt}\setlength{\parskip}{0pt}
\item 全组：比赛 ＋ 演讲
\end{itemize}

\newpage
\sect{5. 里程碑}
{\small
\begin{tabular}{@{}L{2.6cm}L{14.4cm}@{}}
\toprule
\textbf{日期} & \textbf{要做完什么}\tabularnewline
\midrule
9/24（周四） & 原理图定稿 ＋ 车架/PCB 对接约定；设计文件交付\tabularnewline
9/28（周一） & 面包板 demo 交付并开始冒烟；《引脚分配表》冻结；车架模型冻结\tabularnewline
9/30（周三） & PCB 下单（节前窗口）＋ 驱动接口初版\tabularnewline
10/11（周日） & PCB 点亮 ＋ 接线图\tabularnewline
10/13（周二） & 实车 v1 能通电\tabularnewline
10/14（周三） & 驱动接口冻结；第一次上车调试\tabularnewline
10/18（周日） & 循迹初版能跑弯、避障能发现障碍；定版车架\tabularnewline
10/21（周三） & 稳定循迹全赛道；成本台账；展示材料草稿\tabularnewline
10/25（周日） & 全功能冻结：4 障全过 ＋ 精停 ＋ 全套材料\tabularnewline
10/26\textasciitilde{}11/1 & 比赛 ＋ 演讲\tabularnewline
\bottomrule
\end{tabular}

}

\sect{6. 假期安排}
{\small
\begin{tabular}{@{}L{2.6cm}L{2.2cm}L{12.2cm}@{}}
\toprule
\textbf{时间} & \textbf{假期} & \textbf{怎么安排}\tabularnewline
\midrule
9/25\textasciitilde{}9/27 & 中秋节 & 人在校，照常推进；王政淋按原理图搭面包板 demo。\tabularnewline
10/1\textasciitilde{}10/7 & 国庆节 & 实验室和快递基本停。只排软件任务：驱动冒烟、P02、离线算法、文档与素材；硬件一律等 10/8 以后。\tabularnewline
\bottomrule
\end{tabular}

}

\sect{7. 协作纪律}
\begin{enumerate}\setlength{\itemsep}{2pt}
\item 固定顺序：拉取 → 放入 → 提交。开工前双击 \texttt{拉取.bat}，干完双击 \texttt{提交.bat}；提交要组长确认。
\item 掐好节前窗口：PCB 下单 9/30 前（王政淋）、加工材料下单 9/24 前（李昕程）；国庆期间快递与加工停摆。
\item 对接约定定稿后改东西要双方说一声，不能单方面改（《引脚分配表》《车架与 PCB 对接约定》）。
\item 算法谁写的谁上车调，调参不外包。
\item 每次上车测试留日志 + 一句话结论（用 \texttt{tools/run.py} 自动落盘）。
\item 判罚、评分类问题先查 \texttt{参考资料/原文件/} 原文，不确定就在群里问，别自己猜。
\end{enumerate}

\sect{8. 已完成（截至 9/20）}
{\small
\begin{tabular}{@{}L{3.6cm}L{14.2cm}@{}}
\toprule
\textbf{事项} & \textbf{结果}\tabularnewline
\midrule
D01 仓库与协作规范 & 仓库骨架、两个一键脚本、说明书、练习场、各工作目录 README 就绪；5 人全部入仓\tabularnewline
P01 阶段 1 & 一键烧录 ＋ 100Hz 遥测 ＋ 自动存档出图 ＋ 在线调参命令（9/13 跑通）\tabularnewline
引脚与电气设计 v0.2 & \texttt{docs/引脚分配表.md}、\texttt{hardware/原理图草案/}（4 张图 ＋《电气连接规范》＋ 机器可读 \texttt{netlist.json}）\tabularnewline
H01 选型与采购 & BOM 与器件资料索引定稿；万向轮、编码器座、HC-SR04、蓝牙、电源座等已下单；电机确认带编码器；避障定案 HC-SR04 ×2\tabularnewline
M01 建模起步 & 三轮构型整车模型、车架与底板截图、实物照片\tabularnewline
资料与协作 & 参考资料补齐（器件手册、实物照片）；计划全面日期化；\texttt{mech/设计文件/}、\texttt{加工证据/} 与对接约定模板就位\tabularnewline
\bottomrule
\end{tabular}

}

\end{document}
